\section{The Honest Open Question}
\label{sec:open-question}

\begin{itemize}
  \item Neither of the two most naturally testable single-variable accounts
    of the cliff's dissolution survives: the $K/\dstate$ ratio (Act 2)
    is ruled out directly by the flat $d=128$ result, and scalar,
    rank-4-structured anchor-table coherence (Act 3) is ruled out directly
    by the flat dose-response result --- injecting coherence up to and
    beyond $d=64$'s own realized training band still does not reproduce
    the cliff at $d=128$.
  \item \textbf{Two candidates survive, neither confirmed nor excluded by
    any wave reported here.}
  \begin{enumerate}
    \item \textbf{Overlap structure (diffuse vs.\ concentrated).} A
      concentrated (rank-4) injection collapses the table's $n$ rows toward
      a shared low-dimensional subspace; a \emph{diffuse}
      (\texttt{subspace\_rank=48}) injection spreads the same scalar dose
      across a far less structured set of directions. It remains possible
      the cliff tracks \emph{how} coherence is distributed across the
      table, not merely its scalar magnitude. This is a co-primary,
      pre-registered follow-on arm --- design complete, not yet run.
    \item \textbf{Absolute state capacity.} $d_{\mathrm{state}}$ grew
      $4\times$ (64$\to$128) between the cliff's location and its
      disappearance, while $K$ only grew $\sim\!2\times$ (34--46 at $d=64$
      vs.\ 68--92 at $d=128$, holding $K/\dstate$ fixed). It remains
      possible the relevant quantity is raw available state dimensionality
      itself (or slack $\dstate - K$), not any ratio or coherence measure
      --- a hypothesis no wave here can speak to, since $d_{\mathrm{state}}$
      was never varied independently of $K/\dstate$ or coherence.
  \end{enumerate}
  \item Adjudicating between these requires two not-yet-designed waves: a
    diffuse-structure Stage 2 (co-primary with the rank-4 result already
    reported) for candidate 1, and an independent
    $d_{\mathrm{state}}$-vs-$K$/coherence factorial design for candidate 2.
    Neither is registered in detail as of this writing; both are natural,
    concrete next steps, stated here rather than left implicit.
  \item We report this as a genuinely open question, not a hedge: a
    precisely-located transition, its clean dissolution with scale, and
    the ruling-out of the single most natural mechanistic account of that
    dissolution is, on its own, a complete and informative capacity story
    --- even though it does not yet identify what \emph{does} drive the
    transition.
\end{itemize}
